
\begin{proof}[命题~\ref{prop:7-3-10} 的证明]
为证明命题~\ref{prop:7-3-10}，
按 $g=\|(\mathfrak d,\beta)\|$
和 $g^+=\|(\mathfrak d^+,\beta)\|$
归纳构造
$\mathfrak s_{\mathfrak d,k,\beta,\vec P}$
与
$\mathfrak s^+_{\vec a,\mathfrak d^+,k,\beta,\vec P^+}$。
由引理~\ref{lem:7-2-47} 可知，
这样的归纳是可行的。
$\mathfrak s_{\mathfrak d,k,\beta,\vec P}$
的构造与第~\ref{subsec:7-2-5} 小节的论证相似。
下面说明
$\mathfrak s^+_{\vec a,\mathfrak d^+,k,\beta,\vec P^+}$
的构造。
相容条件~\ref{cond:7-3-13}、
\ref{cond:7-3-14} 和
\ref{cond:7-3-17} 归纳地确定
它在
$\partial([0,1]^{|\vec a|})\times
\mathcal M^{\mathrm{main}}_{k+1}
(\beta;\vec P^+)$ 上的限制。
另一方面，(7.3.11.6) 归纳地确定
它在
$[0,1]^{|\vec a|}\times
\partial\bigl(\mathcal M^{\mathrm{main}}_{k+1}
(\beta;\vec P^+)\bigr)$ 上的限制。
因此可由定理 A1.23 将它延拓。
\end{proof}

现在可以定义 $\mathfrak h_{k+1}$。
设 $\vec P^{\,i}\in B_{j_i}C[1]$，
$\vec P^{\,i}=P_{i,1}\otimes\cdots\otimes P_{i,j_i}$
（$i=1,\ldots,\ell+1$），
并置 $j_1+\cdots+j_{\ell+1}=m$，
其中 $\ell\geq0$。
令
\[
\begin{aligned}
\vec a={}&(j_1+1,j_1+j_2+2,\ldots,
j_1+\cdots+j_\ell+\ell),\\
\vec P={}&(P_{1,1},\ldots,P_{1,j_1},
P_{2,1},\ldots,P_{\ell,j_\ell},
P_{\ell+1,1},\ldots,P_{\ell+1,j_{\ell+1}}).
\end{aligned}
\]
这里 $|\vec a|=\ell$，且
$k=\ell+\sum_{i=1}^{\ell+1}j_i$。

\begin{definition}{7.3.18}\label{def:7-3-18}
对 $\ell\geq1$，定义
\begin{equation}
\begin{aligned}
&\mathfrak h_{\ell+1,\beta}
(\vec P^{\,1}\mathbin{\overline\otimes}\cdots
\mathbin{\overline\otimes}\vec P^{\,\ell+1})\\
&\quad=\operatorname{ev}_{0*}\!\left(
\bigl([0,1]^{|\vec a|}\times
\mathcal M^{\mathrm{main},\mathfrak d^+}_{k+1}
(\beta;\vec P^+)\bigr)
^{\mathfrak s^+_{\vec a,\mathfrak d^+,k,\beta,\vec P^+}}
\right).
\end{aligned}
\tag{7.3.19}\label{eq:7-3-19}
\end{equation}
右端模空间的定向见第 8.10 节。
另定义 $\mathfrak h_{1,\beta}=\mathfrak m_\beta$。
\end{definition}

实际上还须稍微修改
$\mathfrak h_{\ell+1,\beta}$；
本节稍后再讨论。

在核对所需性质 (3.3.10) 前，
先看一个特例，以说明上述定义。
考虑 $\mathfrak h_{2,\beta_0}$。
定义它所用的模空间是
$\mathcal M^{\mathrm{main}}_3(\beta_0)$。
回忆作为集合，
$\mathcal M^{\mathrm{main}}_3(\beta_0)=L$；
但其评价映射不是对角嵌入，而是
\[
(\operatorname{ev}_1,\operatorname{ev}_2,
\operatorname{ev}_0)(x)
=\bigl(\varphi_0^{-1}(x),
\varphi_1^{-1}(x),x\bigr),
\]
其中 $\varphi_0,\varphi_1\colon L\to L$
是接近恒等映射的一般微分同胚。
对 $i=0,1$，取短同痕
$\varphi_i^t\colon L\to L$，
使 $\varphi_i^1(x)=x$、
$\varphi_i^0(x)=\varphi_i(x)$。
定义
$\mathcal M^{\mathrm{main}}_3(\beta_0)^+
=[0,1]\times L$，
并赋予评价映射
\[
(\operatorname{ev}_1,\operatorname{ev}_2,
\operatorname{ev}_0)(t,x)
=\bigl((\varphi_0^t)^{-1}(x),
(\varphi_1^t)^{-1}(x),x\bigr).
\]

若 $\vec P^+=(P,L)$、$k=2$、
$\vec a=\{2\}$，则有
\[
\bigl([0,1]^{|\vec a|}\times
\mathcal M^{\mathrm{main},\mathfrak d^+}_3
(\beta;\vec P^+)\bigr)
^{\mathfrak s^+_{\vec a,\mathfrak d^+,k,\beta,\vec P^+}}
=\mathcal M^{\mathrm{main}}_3(\beta_0)^+
\mathop{\times}_{L^2}(P\times L).
\]
参见注~\ref{rem:7-3-15}。
因此
\[
\begin{aligned}
\mathfrak h_{2,\beta_0}
(P\mathbin{\overline\otimes}1)
&=\operatorname{ev}_{0*}\!\left(
\mathcal M^{\mathrm{main}}_3(\beta_0)^+
\mathop{\times}_{L^2}(P\times L)\right)\\
&=\bigcup_{t\in[0,1]}\varphi_0^t(P).
\end{aligned}
\]
这里评价链只有一个时间参数
\translatornote{源书此式在
$\mathcal M^{\mathrm{main}}_3(\beta_0)^+$
前又写 $[0,1]\times$；
但该模空间已定义为 $[0,1]\times L$，
故依紧邻的零集等式省去重复因子。}。
于是
\begin{equation}
\delta\bigl(\mathfrak h_{2,\beta_0}
(P\mathbin{\overline\otimes}1)\bigr)
=P-\varphi_0(P)
=P-\mathfrak m_{2,\beta_0}(P\otimes[L]).
\tag{7.3.20}\label{eq:7-3-20}
\end{equation}
这是 $\mathfrak h$ 所需公式 (3.3.10)
的一部分。
